\scanpage{22}
\mspage{- 3 -}
\margin{n°355}

Le $B$-module $E\otimes_A F$ est réunion des sous-modules $E'\otimes_A F$, $E'$ parcourant
l'ensemble des sous-modules de type fini de $E$. Si donc $\underline P$ appartient à
$\operatorname{Ass}_B(E\otimes_A F)$, il existe un $E'$ de type fini tel que $\underline P\in\operatorname{Ass}_B(E'\otimes_A F)$. D'après
(ii), il existe $\underline p\in\operatorname{Ass}_A(E')$ tel que $\underline P\in\operatorname{Ass}_B(F/\underline pF)$ ; comme
$\operatorname{Ass}_A(E')\subset\operatorname{Ass}_A(E)$, cela achève de démontrer la proposition.

COROLLAIRE.- \uline{Si} $A$ \uline{est noethérien}, \uline{et si} $\underline P\in\operatorname{Ass}_B(E\otimes_A F)$, \uline{on a}
$\pi^{-1}(\underline P)\in\operatorname{Ass}_A(E)$.

Cela résulte des propositions 1 et 2.

\uline{Remarques}.-

1) Supposons $A$ \uline{et} $B$ \uline{noethériens}. Soit $\underline P$ un idéal premier de $B$, et soit
$\underline p=\pi^{-1}(\underline P)$. Soit $S=A-\underline p$, et soit $k(\underline p)=S^{-1}(A/\underline p)$ le corps des fractions
de $A/\underline p$. Puisque $\underline P$ contient $\pi(\underline p)$, il définit un idéal premier de $B/\underline pB$, et
puisqu'il ne rencontre pas $S$, il définit aussi un idéal premier, noté $\underline{P'}$,
de $S^{-1}(B/\underline pB)=k(\underline p)\otimes_A B$. Alors, \uline{pour que} $\underline P\in\operatorname{Ass}_B(E\otimes_A F)$, \uline{il faut et il}
\uline{suffit que} $\underline p\in\operatorname{Ass}_A(E)$ \uline{et} $\underline{P'}\in\operatorname{Ass}_{k(\underline p)\otimes_A B}(k(\underline p)\otimes_A F)$.

Vu la prop.2, tout revient à voir que les conditions
\[
\text{`` }\underline P\in\operatorname{Ass}_B(F/\underline pF)\text{ ''}\quad\text{et}\quad
\text{`` }\underline{P'}\in\operatorname{Ass}_{k(\underline p)\otimes_A B}(k(\underline p)\otimes_A F)\text{ ''}
\]
sont équivalentes, ce qui résulte de la prop.5 du §1 (c'est ici que l'on
utilise l'hypothèse noethérienne sur $B$).

[D'après Grothendieck, l'énoncé ci-dessus est celui que l'on utilise dans
les schémas (le produit tensoriel par $k(\underline p)$ s'interprétant comme le passage
à la ``fibre au-dessus de $\underline p$ '').]

2) La proposition 2 s'applique notamment lorsque $B$ \uline{est plat sur} $A$ et que
l'on prend $F=B$. Le module $E\otimes_A F$ est alors simplement le module $E_{(B)}$
\uline{déduit de} $E$ \uline{par extension des scalaires}.

On recopie ensuite les deux dernières pages de l'Etat 1.

\begin{center}-:-:-:-\end{center}
\margin{21}
